\begin{abstract}
Frequent itemset mining algorithms rely on the anti-monotonicity of support to prune the search space, but this very property ensures that globally rare patterns---those with low overall support---are systematically eliminated.
We observe that many practically important patterns are \emph{globally rare yet locally dense}: they co-occur infrequently across the entire transaction history but exhibit concentrated bursts within narrow temporal windows.
Examples include promotional bundles during flash sales, coordinated anomalous activities during security incidents, and seasonal co-purchases in retail data.
We formalize this observation as the \emph{Rare Dense Pattern} (RDP) mining problem, which requires an itemset to satisfy two simultaneous conditions: (1) global rarity---overall support below a maximum threshold, and (2) local density---at least one sliding-window position where co-occurrence count meets a minimum threshold.
We prove a \emph{Weak Anti-Monotonicity} property: while the combined RDP condition is neither monotone nor anti-monotone, the local density condition alone is anti-monotone, enabling Apriori-style candidate pruning.
Based on this insight, we propose a \emph{Two-Phase Mining} algorithm that first discovers all locally dense patterns using sliding-window Apriori, then filters by global rarity.
We prove that this algorithm is both sound and complete---it discovers exactly the set of all RDPs.
Experiments on synthetic data with embedded ground-truth patterns show 100\% recovery rate across five random seeds, while the rarity filter removes over 80\% of false positives (globally frequent patterns that happen to be locally dense).
The algorithm scales near-linearly to 100K transactions with sub-second runtime.
Our results demonstrate that the RDP framework fills an important gap between rare pattern mining and temporal pattern mining, enabling discovery of actionable patterns invisible to both paradigms individually.
\end{abstract}
